{\small\textbf{A swap, not a squeeze (Diego, 1:15) [15:10–16:25]}}\par\smallskip
Here is what that last row means in euros. This chart is our own approximate calculation.
\begin{itemize}\setlength\itemsep{0pt}
\item In 2020 the systemic-buffer cut released roughly \textbf{six billion} of required capital.
\item The countercyclical buffer then added \textbf{3.3 billion} in 2023 and \textbf{3.4 billion} in 2024.
\item At the same time the buffer for systemically important banks was cut by about \textbf{2.7 billion}.
\item Net, relative to before COVID, the requirement is roughly \textbf{unchanged}, if anything slightly lower.
\end{itemize}
\par\smallskip
DNB said the same: the 2020 plan would be "more or less capital-neutral". In 2023, Governor Knot spoke of "a limited increase in the net capital requirements".
\par\smallskip
Two caveats on our numbers. The cuts apply to the big banks' worldwide assets, while the countercyclical buffer applies to all banks' Dutch exposures, so the bases differ. And measured from 2021 rather than from before COVID, there was a modest increase. We come back to that.
\par\smallskip
So capital was \textbf{moved}, from structural buffers that cannot be released into one that can. The UK did something similar when it set its two-percent neutral rate.
\par\smallskip
But a swap still has a cost, especially in a storm. So, back to our question: was it the wrong moment?
